% META: {"id":"1NSI-TC-EX-007","chapitre":"1NSI-TYPES-CONSTRUITS","type_objet":"exercice","capacites":["1NSI-TYPES-CONSTRUITS-C1"],"parcours":1,"competences":["analyser"],"duree_min":5,"mode_creation":"ex_nihilo","sources_inspiration":[],"status":"approved","origin":"needs_review","status_history":["needs_review","audited","accepted","approved"],"release_acceptance":"RELEASE_OWNER_BATCH_ACCEPTANCE","acceptance_closure_digest":"sha256:9b3ccf9a81c5520fbb7e03b7b2d3e7bf2057a02834b6d908bf3be5f49f3b553f","acceptance_receipt_digest":"sha256:4fad86f0282e0fa4e0419cca1ad6379ae7af5a280c04a4a90880f90a1c044242"}
% BEGIN-TRACE
% match = ("TeamA", "TeamB", 2, 1)
% equipe1, equipe2, score1, score2 = match
% résultat = equipe1 if score1 > score2 else equipe2
% print(résultat)
% print(score1 + score2)
% EXPECTED
% TeamA
% 3
% END-TRACE
\begin{exercice}{1NSI-TC-EX-007}{1}{5}
On modélise un match d'e-sport par un tuple \lstinline{(equipe1, equipe2, score1, score2)}.
\begin{python}
match = ("TeamA", "TeamB", 2, 1)
equipe1, equipe2, score1, score2 = match
resultat = equipe1 if score1 > score2 else equipe2
print(resultat)
print(score1 + score2)
\end{python}
\begin{enumerate}
  \item Donner, sans exécuter le programme, ce qu'affichent les deux \lstinline{print}.
  \item Que contient la variable \lstinline{resultat} si le match est \lstinline{("X", "Y", 0, 3)} ?
\end{enumerate}
\end{exercice}
